\subsection{Lifting of measures}

PROPOSITION 8. --- Let T and X be two topological spaces, \(\pi\) a mapping of T into X.

\begin{enumerate}
\def\labelenumi{\alph{enumi})}
\item
  Let \(\nu\) be a bounded measure on X. In order that there exist a measure \(\mu\) on T such that \(\pi\) is \(\mu\) -proper and \(\pi(\mu) = \nu\) , it is necessary and sufficient that there exist, for every number \(\varepsilon > 0\) , a compact set \(\mathrm{K}_{\varepsilon} \subset \mathrm{T}\) such that the restriction of \(\pi\) to \(\mathrm{K}_{\varepsilon}\) is continuous and \(\nu^{\bullet}(\mathrm{X} - \pi(\mathrm{K}_{\varepsilon})) < \varepsilon\) .
\item
  Suppose that \(\pi\) is injective; let \(\mu\) and \(\mu'\) be two measures on T, such that \(\pi\) is proper for \(\mu\) and \(\mu'\) , and such that \(\pi(\mu) = \pi(\mu')\) . Then \(\mu = \mu'\) .
\end{enumerate}

The condition stated in \(a\) ) is necessary. For, if \(\pi\) is \(\mu\) -proper and \(\pi(\mu) = \nu\) , the relation \(\mu^{\bullet}(1) = \nu^{\bullet}(1) < +\infty\) implies that \(\mu\) is bounded. Prop. 2 of §1, No.~2, applied to the function 1, implies the existence of a compact subset K of T such that \(\mu^{\bullet}(\mathrm{T} - \mathrm{K}) < \varepsilon/2\) . Since \(\pi\) is \(\mu\) -measurable, there exists a compact set \(\mathrm{K}_{\varepsilon} \subset \mathrm{K}\) such that the restriction of \(\pi\) to \(\mathrm{K}_{\varepsilon}\) is continuous, and such that \(\mu^{\bullet}(\mathrm{K} - \mathrm{K}_{\varepsilon}) < \varepsilon/2\) . Then (No.~3, Prop. 4)

\[
\nu^ {\bullet} \big (\mathrm{X} - \pi (\mathrm{K} _ {\varepsilon}) \big) = \mu^ {\bullet} \big (\mathrm{T} - \pi^{-1} \big (\pi (\mathrm{K} _ {\varepsilon}) \big) \big) <   \varepsilon .
\]

To show that the condition is sufficient, we first treat a special case.

Lemma 1. --- Let U and V be two compact spaces, h a continuous mapping of U onto V. The mapping \(\lambda \mapsto h(\lambda)\) of \(\mathcal{M}_{+}(U)\) into \(\mathcal{M}_{+}(V)\) is then surjective.

For, let \(a\) be the linear mapping \(f \mapsto f \circ h\) of \(\mathcal{C}(\mathrm{V})\) into \(\mathcal{C}(\mathrm{U})\) ; since \(h\) is surjective, \(a\) is an isometry of \(\mathcal{C}(\mathrm{V})\) onto a subspace \(\mathrm{H}\) of \(\mathcal{C}(\mathrm{U})\) . Let \(\theta\) be a positive measure on \(\mathrm{V}\) ; then \(\theta \circ a^{-1}\) is a continuous linear form on \(\mathrm{H}\) , which is extendible to a linear form \(\eta\) on \(\mathcal{C}(\mathrm{U})\) with the same norm, by virtue of the Hahn--Banach theorem (TVS, II, §3, No.~2, Cor. 3 of Th. 1); \(\eta\) is then a measure on \(\mathrm{U}\) , and \(\theta(f) = \eta(f \circ h)\) for all \(f \in \mathcal{C}(\mathrm{V})\) , so that \(\theta = h(\eta)\) . Finally, \(\theta(1) = \| \theta \| = \| \eta \|\) , and \(\theta(1) = \eta(1)\) , so that \(\eta\) is positive (Ch. V, §5, No.~5, Prop. 9).

Let us now prove the sufficiency of the condition stated in \(a\) ). The condition implies the existence of a sequence \((\mathbf{K}_n)_{n \geqslant 1}\) of compact subsets of \(T\) , such that the restriction of \(\pi\) to each \(\mathbf{K}_n\) is continuous, and such that, for every \(n\) , \(\nu^{\bullet}(\mathrm{X} - \pi(\mathrm{K}_n)) < 1/n\) . The sequence \((\mathbf{K}_n)\) can be assumed to be increasing. Set \(\mathrm{L}_n = \pi(\mathrm{K}_n)\) and denote by \(\nu_n'\) the measure \(\varphi_{\mathrm{L}_n - \mathrm{L}_{n-1}} \cdot \nu_{\mathrm{L}_n}\) on \(\mathrm{L}_n\) , with the convention \(\mathrm{L}_0 = \emptyset\) .

The restriction \(\pi_{\mathrm{K}_n}\) being continuous, there exists a measure \(\mu_n'\) on \(\mathbf{K}_n\) such that \(\pi_{\mathrm{K}_n}(\mu_n') = \nu_n'\) (Lemma 1). Let \(\mu_n\) be the image of \(\mu_n'\) under the canonical injection of \(\mathbf{K}_n\) into \(\mathbf{T}\) , and let \(g\) be an element of \(\mathcal{F}_{+}(\mathbf{X})\) . Using successively the fact that \(\nu\) is concentrated on \(\bigcup_{n} \mathrm{L}_n\) , Prop. 4 of §1, No.~5; Prop. 2 of §1, No.~2; Prop. 4 of No.~3, and finally Prop. 7 of No.~3, we have

\[
\begin{array}{r l} \nu^ {\bullet} (g) & = \sum_ {n} \nu^ {\bullet} (\varphi_ {\mathrm{L} _ {n} - \mathrm{L} _ {n - 1}} g) = \sum_ {n} \nu_ {n} ^ {\prime   \bullet} (g _ {\mathrm{L} _ {n}}) = \sum_ {n} \mu_ {n} ^ {\prime   \bullet} (g _ {\mathrm{L} _ {n}} \circ \pi_ {\mathrm{K} _ {n}}) \\ & = \sum_ {n} \mu_ {n} ^ {\prime   \bullet} ((g \circ \pi) _ {\mathrm{K} _ {n}}) = \sum_ {n} \mu_ {n} ^ {\bullet} (g \circ \pi). \end{array}
\]

Taking \(g = 1\) in this formula, one sees that the family \((\mu_n)\) is summable and that its sum is a bounded measure \(\mu\) (§1, No.~7, Prop. 7). By Prop. 5 of No.~3, the mapping \(\pi\) is \(\mu_n\) -measurable for all \(n\) , because \(\pi_{\mathrm{K}_n}\) is continuous, hence \(\mu_n'\) -measurable; it follows that \(\pi\) is \(\mu\) -measurable (§1, No.~7, Prop. 8), hence \(\mu\) -proper since \(\mu\) is bounded. The above relations then prove that the measures \(\pi(\mu)\) and \(\nu\) have the same essential upper integral, hence are equal (§1, No.~2, Cor. of Prop. 2).

Finally, let us assume that \(\pi\) is injective, and let us prove \(b\) . Let \(f\) be an element of \(\mathcal{F}_{+}(\mathrm{T})\) ; since \(\pi\) is injective, there exists a function \(g \in \mathcal{F}_{+}(\mathrm{X})\) such that \(f = g \circ \pi\) and, setting \(\nu = \pi(\mu) = \pi(\mu')\) , by Prop. 4 of No.~3 we have

\[
\mu^ {\bullet} (f) = \mu^ {\bullet} (g \circ \pi) = \nu^ {\bullet} (g) = \mu^ {\prime \bullet} (g \circ \pi) = \mu^ {\prime \bullet} (f).
\]

The two measures \(\mu\) and \(\mu'\) thus have the same essential upper integral, which implies their equality (§1, No.~2, Cor. of Prop. 2).

Remark. --- Suppose that \(\pi\) is injective. Let \(\theta\) be a complex measure such that \(\pi\) is \(\theta\) -proper and \(\pi(\theta)=0\) ; then \(\theta=0\) . Indeed, by separating \(\theta\) into its real and imaginary parts, one can reduce to the case that \(\theta\) is real. We then have \(\pi(\theta^{+})=\pi(\theta^{-})\) , therefore \(\theta^{+}=\theta^{-}\) (Prop. 8), and finally \(\theta=0\) .

Here is an important case where condition \(a)\) of Prop. 8 is always satisfied.

PROPOSITION 9. --- Let T be a Souslin space (GT, IX, §6, No.~2, Def. 2), X a Hausdorff space, \(\pi\) a continuous mapping of T onto X, and \(\nu\) a bounded measure on X. Then there exists a bounded measure \(\mu\) on T such that \(\pi(\mu) = \nu\) .

The hypotheses obviously imply that \(\mathbf{X}\) is a Souslin space.

Let us consider the set function \(c: \mathrm{A} \mapsto \nu^{\bullet}(\pi(\mathrm{A}))\) on \(\mathfrak{P}(\mathrm{T})\) . The relation \(\mathrm{A} \subset \mathrm{B}\) implies \(c(\mathrm{A}) \leqslant c(\mathrm{B})\) ; if \((\mathrm{A}_n)\) is an increasing sequence of subsets of \(\mathrm{T}\) , and if \(\mathrm{A} = \bigcup_{n \in \mathbf{N}} \mathrm{A}_n\) , then \(c(\mathrm{A}) = \sup_n c(\mathrm{A}_n)\) from the fact that \(\nu^{\bullet}\) is an encumbrance. Finally, let \(\mathrm{A} \subset \mathrm{T}\) and let \(\varepsilon\) be a number \(> 0\) ; choose an open subset \(\mathrm{G}\) of \(\mathrm{X}\) containing \(\pi(\mathrm{A})\) , such that \(\nu^{\bullet}(\mathrm{G}) \leqslant \nu^{\bullet}(\pi(\mathrm{A})) + \varepsilon\) (§1, No.~9, Prop. 13); the open subset \(\mathrm{H} = \overline{\pi}^{-1}(\mathrm{G})\) of \(\mathrm{T}\) contains \(\mathrm{A}\) , and \(c(\mathrm{H}) \leqslant c(\mathrm{A}) + \varepsilon\) . The function \(c\) is therefore a right-continuous capacity on \(\mathrm{T} (\mathrm{TG},\mathrm{IX},\S 6,\mathrm{No.~10},\mathrm{Def.~9})^{(1)}\) and the theorem on capacitability (loc. cit., Th. 6) implies the equality \(c(\mathrm{T}) = \sup_{\mathrm{K}}c(\mathrm{K})\) , where \(K\) runs over the set of compact subsets of \(T\) . Prop. 8 then implies the existence of the desired measure \(\mu\) .

